\input{../tex-inputs/latex-leseansicht-vorspann}

\section[Arthur Schnitzler an Gustav Schwarzkopf, 6. 6. 1898]{L04505 Arthur Schnitzler an Gustav Schwarzkopf, 6. 6. 1898}
\nopagebreak\mylabel{L04505v}
\rehead{ }\normalsize\beginnumbering\briefempfaengerindex{Schwarzkopf, Gustav@\textsc{Schwarzkopf, Gustav}!zzzSchnitzler, Arthur@\emph{von Arthur Schnitzler}!1898-06-061@{6. 6. 1898}|(be}
\toendnotes[C]{\smallbreak\pagebreak[2]}
\correspDesc{Versand  durch Arthur Schnitzler am 6. 6. 1898 in Friesach [Sankt Veit an der Glan]
\newline{}Übermittlung  am 6. 6. 1898 in Friesach [Sankt Veit an der Glan]
\newline{}Zustellung  am 7. 6. 1898 in I., Innere Stadt
\newline{}Erhalt  durch Gustav Schwarzkopf im Zeitraum [7. 6. 1898 – 11. 6. 1898?] in Wien}\toendnotes[C]{\smallbreak}
\Standort{DLA, A:Schnitzler, HS.1985.1.1897.}
\physDesc{Bildpostkarte, 75 Zeichen
\newline{}Handschrift: Bleistift, deutsche Kurrent
\newline{}Versand: 1) Stempel: »\nobreak{}\oindex{Friesach [Sankt Veit an der Glan]@\textbf{Friesach [Sankt Veit an der Glan]}|pwk}Friesach, 6 6 98\nobreak{}«.   2) Stempel: »\nobreak{}\oindex{I., Innere Stadt@\textbf{I., Innere Stadt}|pwk}Wien 1/1 1, 7\textcolor{gray}{.} 6. 98, 9–10½ V., bestellt\nobreak{}«. }
\pstart
           \noindent{}\raggedleft{}{\pb}\textcolor{gray}{\textbf{Friesach\oindex{Friesach [Sankt Veit an der Glan]@\textbf{Friesach [Sankt Veit an der Glan]}|pw}
               .}}\pend
           
\pstart
           \centering{}\textcolor{gray}{\textbf{Platz\oindex{Hauptplatz@\textbf{Hauptplatz}|pw}.}}\pend
           
\pstart
           \centering{}\textcolor{gray}{\textbf{Mathias Hayd\pwindex{Hayd, Mathias @\textsc{Hayd, Mathias}|pw}, Friesach\oindex{Friesach [Sankt Veit an der Glan]@\textbf{Friesach [Sankt Veit an der Glan]}|pw}.}}\pend
           \pstart{}{\pb}\textsc{Herrn Gustav Schwarzkopf}\pend{}\pstart{}Wien\oindex{Wien@\textbf{Wien}|pw}\pend{}\pstart{}\textsc{I. Tiefer Graben 23}\oindex{Wien@\textbf{Wien}!I., Innere Stadt@\textbf{I., Innere Stadt}!Tiefer Graben 23@\textbf{Tiefer Graben 23}|pw}.\pend{}{\bigskip}\vspace{1em}
\pstart
           \noindent{}Herzlichen Gruß!\pend
           \pstart Ihr \spacefill\mbox{ArthSchn}\pend{}\selectlanguage{ngerman}\endnumbering\briefempfaengerindex{Schwarzkopf, Gustav@\textsc{Schwarzkopf, Gustav}!zzzSchnitzler, Arthur@\emph{von Arthur Schnitzler}!1898-06-061@{6. 6. 1898}|)be}\mylabel{L04505h}
\begin{anhang}
\end{anhang}\newcommand{\dateiname}{L04505}\newcommand{\titel}{Arthur Schnitzler an Gustav Schwarzkopf, 6. 6. 1898}\newcommand{\editorInnen}{Herausgegeben von Jahnke, SelmaMüller, Martin Anton}\input{../tex-inputs/latex-leseansicht-abspann}
